% Detailed derivation gathered at the end of the owning structure.

\begin{derivation}[从\eqnrefs{eq:fock-oneparticle-expansion-dp68}到\eqnrefs{eq:fock-oneparticle-schrodinger-dp68}]
一粒子波函数不是额外附加到场论上的对象，而是 Fock 态在局域一粒子基底上的坐标。先从正文的态展开
\begin{equation*}
|\Psi_1(t)\rangle
=\int\dd^3r\,\psi_{\rm 1p}(\mathbf r,t)
\hat\Psi^\dagger(\mathbf r)|0\rangle
\end{equation*}
出发。自由 Schr\"odinger 场 Hamiltonian 为
\begin{equation*}
\hat H
=\int\dd^3x\,
\hat\Psi^\dagger(\mathbf x)
\left(-\frac{\hbar^2\nabla_{\mathbf x}^2}{2m}\right)
\hat\Psi(\mathbf x).
\end{equation*}
利用等时对易关系，先求产生一个局域粒子时能量算符怎样作用：
\begin{align*}
[\hat H,\hat\Psi^\dagger(\mathbf r)]
&=\int\dd^3x\,
\hat\Psi^\dagger(\mathbf x)
\left(-\frac{\hbar^2\nabla_{\mathbf x}^2}{2m}\right)
[\hat\Psi(\mathbf x),\hat\Psi^\dagger(\mathbf r)]\\
&=\int\dd^3x\,
\hat\Psi^\dagger(\mathbf x)
\left(-\frac{\hbar^2\nabla_{\mathbf x}^2}{2m}\right)
\delta^{(3)}(\mathbf x-\mathbf r)\\
&=-\frac{\hbar^2}{2m}\nabla_{\mathbf r}^2
\hat\Psi^\dagger(\mathbf r).
\end{align*}
因为真空满足 \(\hat H|0\rangle=0\)，于是
\begin{align*}
\hat H|\Psi_1\rangle
&=\int\dd^3r\,\psi_{\rm 1p}(\mathbf r)
[\hat H,\hat\Psi^\dagger(\mathbf r)]|0\rangle\\
&=\int\dd^3r\,
\left[-\frac{\hbar^2}{2m}\nabla^2\psi_{\rm 1p}(\mathbf r)\right]
\hat\Psi^\dagger(\mathbf r)|0\rangle,
\end{align*}
其中第二行对空间变量分部积分两次，并假设波函数在边界处足够快衰减。

连续基底中的系数比较由场算符投影严格实现。对任意 \(\mathbf x\)，左乘 \(\langle0|\hat\Psi(\mathbf x)\)，由
\begin{equation*}
\langle0|\hat\Psi(\mathbf x)
\hat\Psi^\dagger(\mathbf r)|0\rangle
=\delta^{(3)}(\mathbf x-\mathbf r)
\end{equation*}
可得
\begin{align*}
\langle0|\hat\Psi(\mathbf x)
\ii\hbar\partial_t|\Psi_1\rangle
&=\ii\hbar\partial_t\psi_{\rm 1p}(\mathbf x,t),\\
\langle0|\hat\Psi(\mathbf x)
\hat H|\Psi_1\rangle
&=-\frac{\hbar^2}{2m}\nabla_{\mathbf x}^2
\psi_{\rm 1p}(\mathbf x,t).
\end{align*}
因此 Fock 空间 Schr\"odinger 方程 \(\ii\hbar\partial_t|\Psi_1\rangle=\hat H|\Psi_1\rangle\) 在每个 \(\mathbf x\) 上严格投影成
\begin{equation*}
\ii\hbar\partial_t\psi_{\rm 1p}(\mathbf x,t)
=-\frac{\hbar^2}{2m}\nabla_{\mathbf x}^2\psi_{\rm 1p}(\mathbf x,t),
\end{equation*}
即正文\eqnrefs{eq:fock-oneparticle-schrodinger-dp68}。这里真正建立的是“固定粒子数 Fock 扇区”与普通单粒子波函数之间的坐标同构，而不是重新假设一条量子力学方程。
\end{derivation}
